\section{Parallel Planning Approach}
\label{sec:parallel_planning_approach}
\hspace{\parindent}
The proposed architecture extends the original TEA*-based planner by executing several
independent TEA* instances in parallel. As illustrated in
Figure~\ref{fig:parallel_planning_framework}, the Supervisor sends a planning request to the
Path Planner containing the current robot positions, their corresponding goals, and the maximum
time available for planning, denoted by
\(I=\{(p_i,g_i)\}_{i=1}^{N}, T_{\max}\).

After receiving the request, the Path Planner initializes the execution and spawns a set of worker
threads. Each thread retrieves a heuristic from a shared heuristic pool, uses it to compute a
priority permutation, and then runs TEA* independently using that ordering. Each execution
therefore produces a candidate plan associated with one priority ordering.

\begin{figure}[H]
    \centering
    \includegraphics[width=1.0\linewidth]{parallel_planning_framework.png}
    \caption{Overview of the proposed parallel planning framework.}
    \label{fig:parallel_planning_framework}
\end{figure}


For simplicity, Figure~\ref{fig:parallel_planning_framework} represents the case in which the
number of available heuristics is equal to the number of worker threads. In that simplified case,
each thread retrieves one heuristic and produces one candidate plan. In the general case, however, these numbers do not need to be equal, since a thread may evaluate several heuristics during the same planning cycle as long as planning time is still available and heuristics remain to be tested.

During the parallel execution stage, each worker thread generates one or more candidate plans, which are collected in a set \(P\) and subsequently evaluated using the selected cost function \(J\). Once all candidate plans have been evaluated, the plan with the lowest cost is returned to the Supervisor. In this work, \(J(P_i)\) is defined as the sum of the individual costs of all robot paths in candidate plan \(P_i\). Candidate plans that do not include valid paths for all robots are assigned an \(\infty\) cost, since a plan is considered to have failed whenever at least one robot is unable to obtain a valid path.

The following subsections describe these components in more detail, namely the thread implementation details, the heuristic dispatching strategy, the termination mechanisms, and the handling of redundant heuristics.


\subsection{Thread Implementation Details}
\label{subsec:thread_implementation_details}
bla bla bla 

\subsection{Heuristic Dispatching Mechanism}
\label{subsec:heuristic_dispatching}
\hspace{\parindent}
The dispatching mechanism was designed to avoid assigning highly redundant heuristics to the
worker threads at the same time. Since the heuristic set is organized into different families, the
objective is to promote a balanced use of these families during the parallel execution. This increases
the diversity of the generated priority orderings and reduces the probability that several threads
evaluate very similar planning alternatives.

To achieve this, the available heuristics are stored in a shared priority queue. Each heuristic is
associated with a heuristic group, and each group has a counter indicating how many heuristics from
that group have already been selected during the current planning cycle. The priority queue is
ordered according to these counters, so that heuristics belonging to the least-used groups are given
higher priority.

Whenever a worker thread needs a new heuristic, it accesses the shared queue and retrieves a
heuristic from the group with the lowest usage counter. After the heuristic is selected, the counter
of the corresponding group is incremented, and the queue is reordered accordingly. The selected
heuristic is also removed from the queue, preventing two threads from evaluating the exact same
heuristic during the same planning cycle.

This mechanism encourages the simultaneous evaluation of heuristics from different families. For
example, instead of assigning several conflict-based heuristics to different threads at the same time,
the dispatcher tends to distribute the available computational resources across distance-based,
topological, surroundings-based, and conflict-based strategies. In this way, the parallel framework
does not merely execute several TEA* instances concurrently; it also attempts to make those
executions complementary.

Since the heuristic pool is shared by all worker threads and is dynamically modified during
execution, access to this structure must be synchronized. For this reason, a mutex is used to protect
the critical section in which a thread retrieves a heuristic, updates the corresponding group counter,
removes the selected heuristic from the queue, and restores the queue ordering. This ensures that
only one thread modifies the heuristic pool at a time, preventing race conditions and maintaining a
consistent dispatching state throughout the planning cycle.

\subsection{Timeout Handling and Controlled Termination}
\label{subsec:timeout_controlled_termination}

\subsection{Handling Redundant Heuristics}
\label{subsec:handling_redundant_heuristics}





\newpage
\section{Summary}
\label{sec:parallel_summary}